\begin{tikzpicture}
    \begin{umlpackage}[]{WebRTC}
        \umlenum[x=0,y=0]{Description}{
            + Answer(RTCSessionDescription) \\
            + Offer(RTCSessionDescription) \\
            + None
        }{}

        \umlclass[x=0,y=-6,type=ThreadSafeShared]{WebRTCManager}{
            + peer\_connection : RTCPeerConnection \\
            + ice : Vector<RTCIceCandidate> \\
            + channel : RTCDataChannel \\
            + session : Optional<Session> \\
            + description : Description
        }{
            \umlvirt{+ init(servers : Vector<RTCIceServer>) : WebRTCManager} \\
            \umlvirt{+ create\_offer() : RTCSessionDescription} \\
            \umlvirt{+ create\_answer() : RTCSessionDescription} \\
            \umlvirt{+ connect(offer: String) : RTCSessionDescription} \\
            \umlvirt{+ create\_channel() : RTCDataChannel} \\
            \umlvirt{+ send(message : String) : << raises Error >>}
        }{}

        \umlclass[x=9,y=-6]{Session}{
            - session\_keys : SessionKeys \\
            - sending\_ratchet : DoubleRatchet \\
            - receiving\_chains : ChainStore
        }{
            \umlvirt{+ encrypt(message : Vector<Byte>) : Vector<Byte>} \\
            \umlvirt{+ decrypt(message : Vector<Byte>) : Vector<Byte>}
        }{}

        \umlaggreg[arg=hasDescription, mult=1, pos=0.5]{Description}{WebRTCManager}
        \umlaggreg[mult=1, pos=0.5]{WebRTCManager}{Session}
    \end{umlpackage}
\end{tikzpicture}

\section{WebRTC Manager Cardinality}

A single instance of the \texttt{WebRTCManager} MUST be dedicated to exactly
one WebRTC session. Implementations MUST NOT reuse a \texttt{WebRTCManager}
instance across multiple concurrent or sequential sessions to ensure state
isolation.

\section{Peer-to-Peer Handshake Sequence}

The establishment of a secure communication channel follows a strict phased
approach:

\begin{enumerate}
    \item \textbf{WebRTC Signaling}: Peers MUST complete the standard signaling
        exchange (Offer, Answer, and Pranswer/Acknowledgment).

    \item \textbf{Cryptographic Material Exchange}: Immediately following the
        WebRTC connection establishment, both peers MUST exchange their
        cryptographic identity material. This includes:
        \begin{itemize}
            \item The Identity Public Key (\(IK\)).
            \item A Signed Pre-key (\(SPK\)).
            \item A set of One-Time Pre-keys (\(OTK\)).
        \end{itemize}

        \item \textbf{Session Initialization}: Upon successful exchange of
        these keys, peers MUST execute the X3DH (Extended Triple
        Diffie-Hellman) protocol to derive a shared secret. The Session object
        MUST be initialized only after the successful derivation of this
        secret.
\end{enumerate}

\section{Transport Requirements}

The underlying transport layer MUST operate over DTLS (Datagram Transport Layer
Security), as specified in Section 4.2 (Transport Layer). All application data
MUST be encapsulated within the DTLS tunnel to guarantee confidentiality and
integrity.